\documentclass[../../main]{subfiles}

\begin{document}

\section{Funções Características}
\label{sec:matematica:funcoes:funcoes-caracteristicas}

\subsection{Definição}

\begin{defn}[Função Característica]
  \label{def:matematica:funcoes:caracteristicas:funcao-caracteristica}

  Seja \(X\) um conjunto e seja

  \[
    A\subseteq X.
  \]

  A função característica de \(A\) é a função

  \[
    \chi_A:X\to\{0,1\}
  \]

  definida por

  \[
    \chi_A(x)
    =
    \begin{cases}
      1, & \text{se } x\in A,\\
      0, & \text{se } x\notin A.
    \end{cases}
  \]

\end{defn}

\begin{obs}

  Alguns autores utilizam a notação

  \[
    \mathbf 1_A
  \]

  em lugar de

  \[
    \chi_A.
  \]

\end{obs}

\begin{ex}

  Seja

  \[
    X=\mathbb Z
  \]

  e

  \[
    P
    =
    \{
      n\in\mathbb Z
      \mid
      n \text{ é par}
    \}.
  \]

  Então

  \[
    \chi_P(n)
    =
    \begin{cases}
      1, & \text{se } n \text{ é par},\\
      0, & \text{caso contrário}.
    \end{cases}
  \]

\end{ex}

\subsection{Correspondência com Subconjuntos}

\begin{prop}
  \label{prop:matematica:funcoes:caracteristicas:recuperacao}

  Seja

  \[
    A\subseteq X.
  \]

  Então

  \[
    A
    =
    \chi_A^{-1}[\{1\}].
  \]

\end{prop}

\begin{proof}

  Seja \(x\in X\).

  Temos

  \[
    x\in\chi_A^{-1}[\{1\}]
  \]

  se, e somente se,

  \[
    \chi_A(x)=1.
  \]

  Pela definição de função característica,

  \[
    \chi_A(x)=1
  \]

  se, e somente se,

  \[
    x\in A.
  \]

\end{proof}

\begin{prop}
  \label{prop:matematica:funcoes:caracteristicas:bijeccao}

  Seja \(X\) um conjunto.

  A aplicação

  \[
    \Phi
    :
    \mathcal P(X)
    \to
    \{0,1\}^X
  \]

  definida por

  \[
    \Phi(A)=\chi_A
  \]

  é uma bijeção.

\end{prop}

\begin{proof}

  Injetividade:

  Se

  \[
    \chi_A=\chi_B,
  \]

  então para todo \(x\in X\),

  \[
    x\in A
    \iff
    \chi_A(x)=1
    \iff
    \chi_B(x)=1
    \iff
    x\in B.
  \]

  Logo

  \[
    A=B.
  \]

  Sobrejetividade:

  Seja

  \[
    f:X\to\{0,1\}.
  \]

  Defina

  \[
    A=f^{-1}[\{1\}].
  \]

  Então

  \[
    f=\chi_A.
  \]

\end{proof}

\begin{cor}

  Para todo conjunto \(X\),

  \[
    |\mathcal P(X)|
    =
    |\{0,1\}^X|.
  \]

\end{cor}

\subsection{Caracterização da Igualdade}

\begin{prop}
  \label{prop:matematica:funcoes:caracteristicas:igualdade}

  Sejam

  \[
    A,B\subseteq X.
  \]

  Então

  \[
    A=B
    \iff
    \chi_A=\chi_B.
  \]

\end{prop}

\begin{proof}

  Consequência imediata da injetividade de \(\Phi\).

\end{proof}

\subsection{Operações Conjuntistas}

\begin{prop}[Complemento]
  \label{prop:matematica:funcoes:caracteristicas:complemento}

  Para todo

  \[
    A\subseteq X,
  \]

  vale

  \[
    \chi_{X\setminus A}
    =
    1-\chi_A.
  \]

\end{prop}

\begin{proof}

  Verificação ponto a ponto.

\end{proof}

\begin{prop}[Interseção]
  \label{prop:matematica:funcoes:caracteristicas:intersecao}

  Para quaisquer

  \[
    A,B\subseteq X,
  \]

  vale

  \[
    \chi_{A\cap B}
    =
    \chi_A\chi_B.
  \]

\end{prop}

\begin{proof}

  Verificação ponto a ponto.

\end{proof}

\begin{prop}[União]
  \label{prop:matematica:funcoes:caracteristicas:uniao}

  Para quaisquer

  \[
    A,B\subseteq X,
  \]

  vale

  \[
    \chi_{A\cup B}
    =
    \chi_A+\chi_B-\chi_A\chi_B.
  \]

\end{prop}

\begin{proof}

  Verificação ponto a ponto.

\end{proof}

\begin{prop}[Diferença]
  \label{prop:matematica:funcoes:caracteristicas:diferenca}

  Para quaisquer

  \[
    A,B\subseteq X,
  \]

  vale

  \[
    \chi_{A\setminus B}
    =
    \chi_A(1-\chi_B).
  \]

\end{prop}

\begin{proof}

  Verificação ponto a ponto.

\end{proof}

\subsection{Famílias de Conjuntos}

\begin{prop}
  \label{prop:matematica:funcoes:caracteristicas:intersecao-familia}

  Seja

  \[
    (A_i)_{i\in I}
  \]

  uma família de subconjuntos de \(X\).

  Então

  \[
    \chi_{\bigcap_{i\in I}A_i}
    =
    \prod_{i\in I}\chi_{A_i},
  \]

  sempre que o produto faça sentido.

\end{prop}

\begin{proof}

  Um elemento pertence à interseção precisamente quando pertence a
  todos os conjuntos da família.

\end{proof}

\begin{prop}
  \label{prop:matematica:funcoes:caracteristicas:uniao-familia}

  Seja

  \[
    (A_i)_{i\in I}
  \]

  uma família de subconjuntos de \(X\).

  Então

  \[
    \chi_{\bigcup_{i\in I}A_i}(x)
    =
    \sup_{i\in I}
    \chi_{A_i}(x).
  \]

\end{prop}

\begin{proof}

  Segue da definição de união.

\end{proof}

\subsection{Partições}

\begin{prop}
  \label{prop:matematica:funcoes:caracteristicas:particao}

  Seja

  \[
    (A_i)_{i\in I}
  \]

  uma partição de \(X\).

  Então

  \[
    \sum_{i\in I}
    \chi_{A_i}(x)
    =
    1
  \]

  para todo

  \[
    x\in X.
  \]

\end{prop}

\begin{proof}

  Cada elemento de \(X\) pertence a exatamente um conjunto da
  partição.

\end{proof}

\end{document}
